\subsection{Structure of locally compact fields}

The completions R and \(Q_{p}\) of the field Q with respect to the non-improper absolute values on Q (p any prime) are locally compact. On the other hand, for every power \(q = p^{f}\) of a prime number p, the field \(\mathbf{F}_{q}((\mathrm{T}))\) of formal power series over the finite field \(F_{q}\) with the valuation defined in §3, no. 4, Example 3 is locally compact: for the maximal ideal of the valuation ring \(F_{q}[[T]]\) is generated by T; we know that this ring is complete with the (T)-adic topology (Chapter III, §2, no. 6, Proposition 6) and, as the residue field \(F_{q}\) is finite, Proposition 2 of §5, no. 1 proves our assertion. Conversely:

THEOREM 1. Let K be a (not necessarily commutative) locally compact field which is not discrete.

\begin{enumerate}
\def\labelenumi{(\roman{enumi})}
\item
  If \(\mathbf{K}\) is of characteristic 0 and \(\mathrm{mod}_{\mathbf{K}}\) is not an ultrametric absolute value then \(\mathbf{K}\) is isomorphic to one of the fields \(\mathbf{R}\) , \(\mathbf{C}\) or \(\mathbf{H}\) .
\item
  If \(\mathbf{K}\) is of characteristic 0 and \(\mathrm{mod}_{\mathbf{K}}\) is an ultrametric absolute value, \(\mathbf{K}\) is an algebra of finite rank over a \(p\) -adic field \(\mathbf{Q}_p\) .
\item
  If \(\mathbf{K}\) is of characteristic \(p \neq 0\) , it is isomorphic to a field with centre a field of formal power series \(\mathbf{F}_q((\mathrm{T}))\) (where \(q\) is a power of \(p\) ) and of finite rank over its centre.
\end{enumerate}

\begin{enumerate}
\def\labelenumi{(\roman{enumi})}
\item
  It follows from Ostrowski's Theorem (§ 6, no. 4, Theorem 2) that K is a topological field isomorphic to an everywhere dense subfield of R, C or H and, as K is complete, it is isomorphic to R, C or H.
\item
  Let A be the ring of the absolute value \(mod_{K}\) and m its maximal ideal. We know that A/m is a finite field (§ 5, no. 1, Proposition 2) and hence the absolute value induced by \(mod_{K}\) on Q has a finite residue field, which is only possible if it is equivalent to a \(p\) -adic absolute value (§ 6, no. 3, Proposition 4); the closure of \(\mathbf{Q}\) in \(\mathbf{K}\) is therefore isomorphic to \(\mathbf{Q}_p\) and is contained in the centre of \(\mathbf{K}\) since the latter is closed in \(\mathbf{K}\) ; we conclude using Proposition 2 of no. 1.
\item
  The second assertion follows from the first and the Corollary to Proposition 2 of no. 2. To show the first assertion, note that \(\mathrm{mod}_{\mathbf{K}}\) is necessarily an ultrametric absolute value (§ 6, no. 2, Corollary to Proposition 3); in the notation of the proof of Proposition 3 of no. 2, the centre \(Z\) of \(K\) consists of the elements which commute with both \(s\) and \(u\) ; but by virtue of (3),
\end{enumerate}

\[
u ^ {- q} s u ^ {q} = s ^ {q p ^ {r}} = s,
\]

so that \(u^q \in \mathbf{Z}\) and we conclude that \(\mathbf{Z}\) is not discrete. As \(\mathbf{Z}\) is locally compact, we may confine our attention to the case where \(\mathbf{K}\) is commutative. The sub- \(\mathbf{F}_p\) -algebra \(\mathbf{F}_p[s]\) in \(\mathbf{K}\) is then a finite field since \(s^{q-1} = 1\) and obviously \(y^q = y\) for every element of this field, which is therefore identical with \(S\) and isomorphic to \(\mathbf{F}_q\) since \(S \subset \mathbf{F}_p[s]\) has \(q\) elements. Since the sum of two elements of \(S\) is in \(S\) , the mapping which maps each formal power series \(\sum_{i=0}^{\infty} s_i T^i \in \mathbf{F}_q[[T]]\) to the element \(\sum_{i=0}^{\infty}s_{i}u^{i}\) is a bijective homomorphism of the ring \(F_{q}[[T]]\) onto the ring A, whence immediately the conclusion.

COROLLARY 1. Every locally compact field which is not discrete is of finite rank over its centre.

COROLLARY 2. Every locally compact field is connected or totally disconnected; if it is connected, it is isomorphic to R, C or H.

If the topology on a field K is defined by an ultrametric absolute value, K is totally disconnected with this topology.

Remark. Let \(s\) be an integer \(>0\) ; the subfield \(\mathbf{F}_{q}((\mathrm{T}^{s})) = \mathbf{L}\) of \(\mathbf{K} = \mathbf{F}_{q}((\mathrm{T}))\) is closed in \(\mathbf{K}\) and \(e(\mathrm{K/L}) = s\) and \(f(\mathrm{K/L}) = 1\) . It is therefore seen that there are closed subfields \(\mathbf{L}\) of \(\mathbf{K}\) which are not discrete and such that \(e(\mathrm{K/L})\) (and a fortiori the degree {[}K: L{]}) is arbitrarily large (contrary to what happens for locally compact fields of characteristic 0, where every locally compact subfield \(\mathbf{L}\) of such a field \(\mathbf{K}\) necessarily contains \(\mathbf{R}\) or \(\mathbf{Q}_{p}\) and therefore {[}K: L{]} is bounded).

